\begin{figure}[h]
\centering
\begin{verbatim}
+-----------------------------------------------------------+
|                   Simulation Architecture                 |
+-----------------------------------------------------------+
| 1. Load Schema ------------------------------------------ |
|    - field definitions (phi, V, R)                        |
|    - operator coefficients                                |
|    - boundary conditions                                  |
|    - nonlinearities and thresholds                        |
+-----------------------------------------------------------+
| 2. Initialize Fields -------------------------------------|
|    phi(x,0), V(x,0), R(x,0)                               |
|    from config: gradient, rotation, uniform, noise        |
+-----------------------------------------------------------+
| 3. Operator Pipeline -------------------------------------|
|    For each timestep:                                     |
|      a) Diffusion:        D_phi ∇²phi, D_V ∇²V, D_R ∇²R   |
|      b) Alignment:        α * (targets - fields)          |
|      c) Coupling:         β * cross-field interactions    |
|      d) Activation:       γ * nonlinear terms             |
|      e) Stabilization:    λ * damping                     |
+-----------------------------------------------------------+
| 4. Coherence Metrics -------------------------------------|
|    C(phi, V) = (∇phi · V) / (||∇phi|| ||V|| + ε)          |
|    R dynamics include logistic + coherence gain           |
+-----------------------------------------------------------+
| 5. Output -------------------------------------------------|
|    - field snapshots                                      |
|    - coherence curves                                     |
|    - resonance envelope evolution                         |
+-----------------------------------------------------------+
\end{verbatim}
\caption{Text-based simulation architecture diagram showing the full
operator pipeline and data flow.}
\label{fig:sim-architecture}
\end{figure}
